% Reserved D.67 candidate. Unregistered pending independent proof audit.
\subsection{Generic cusp rigidity in the full free pro-five group}

\begin{statement}{D.67}{Proved}{A comeagre set of cusp abelian classes has a unique full pro-five section}
Let \(F_5=F^{(5)}\) be the actual geometric free pro-five group of
the tripod over \(\mathbb Q\), with the unit tangential arithmetic
action at zero. Give
\({\cal K}=H^1(G_{\mathbb Q},\mathbb Z_5(1))\) the intrinsic
uniform five-adic topology of (D.66.2). There is a dense
\(G_\delta\) subset \({\cal G}_{\rm full}\subset{\cal K}\)
such that, for every \(a\in{\cal G}_{\rm full}\), the fibre
\[
 H^1_{\mathrm{cts}}(G_{\mathbb Q},F_5)
       \longrightarrow H^1(G_{\mathbb Q},F_5^{\rm ab})
                       ={\cal K}^2
 \qquad\text{over }(a,0)                               \tag{D.67.1}
\]
is the singleton consisting of the completed cusp section
at zero. This section is supplied by the actual full
profinite cusp packet. The analogous assertion holds in
each of the other two cusp packets.

The set may be chosen so that every parameter is infinitely
ramified and also belongs to the generic singleton
metabelian-fibre locus of D.66. Thus both the full free
pro-five fibre and the metabelian fibre are then singletons.
Every compatible tower of nilpotent section classes
above such an abelian pair is the tower of this cusp class.
No assertion is made that individual lifts to arbitrarily
high finite nilpotent levels admit compatible choices,
or that a first obstructed finite level has been determined.
\end{statement}

\begin{proof}
Write \(P=F_5\), let \(H=P'\) be its closed commutator
subgroup, and let \(\sigma_0:G_{\mathbb Q}\to\operatorname{Aut}(P)\)
be the actual unit tangential action. In particular
\(\sigma_0(g)(x)=x^{\chi_5(g)}\). Use the fixed continuous
Kummer cocycles of D.66 for all parameters. The cusp
section with parameter \(a\) gives the genuine action
\[
 \sigma_a(g)=\operatorname{Int}(x^{a(g)})\sigma_0(g)
                  \quad\text{on }P\text{ and on }H.     \tag{D.67.2}
\]
The cocycle identity for \(a\) proves that (D.67.2) is
an action. A section class with abelian class \((a,0)\)
can be conjugated so that its abelian cocycle agrees
exactly with that of this cusp section. Its difference
is then an actual nonabelian cocycle
\[
                           c:G_{\mathbb Q}\longrightarrow H
                                                        \tag{D.67.3}
\]
for \(\sigma_a\). Conversely every such cocycle gives
a section in that fibre. A conjugator between two
such exact representatives has abelian image fixed
by \(G_{\mathbb Q}\), which is zero in
\(\mathbb Z_5(1)^2\). Hence the fibre is precisely
the nonabelian set \(H^1(G_{\mathbb Q},H;\sigma_a)\).
We shall show that this pointed set is trivial on a
dense \(G_\delta\) of parameters.

\paragraph{The exact local consequence of good reduction.}
Fix a rational prime \(q\ne5\). The tripod with its
three marked boundary points has good reduction over
\(\mathbb Z_q\), and the honest unit tangential action
on \(P\) is unramified. Let \(\Pi_q=P\rtimes G_{\mathbb Q_q}\),
and let \(\Pi_q^{\rm sp}\) be its geometrically pro-five
good-reduction specialization group over
\(G_{\mathbb F_q}\). The comparison diagram is
\[
\begin{array}{ccccccccc}
1&\longrightarrow&P&\longrightarrow&\Pi_q&\longrightarrow&G_{\mathbb Q_q}&\longrightarrow&1\\
 &&\Vert&&\downarrow&&\downarrow\\
1&\longrightarrow&P&\longrightarrow&\Pi_q^{\rm sp}&\longrightarrow&G_{\mathbb F_q}&\longrightarrow&1.
\end{array}                                               \tag{D.67.4}
\]
It is cartesian. The geometric isomorphism and this
description are exactly the prime-set \(\Sigma=\{5\}\)
case of \cite[Definition A.3 and Remark A.3.1, p.~38]{DE_Hoshi2012}.
The residue characteristic is \(q\), and is excluded
from \(\Sigma\), as that statement requires.

For every actual local section, its inertia image
in \(\Pi_q^{\rm sp}\) is either trivial or a
nontrivial subgroup of a conjugate cusp inertia
\(\mathbb Z_5(1)\). This is
\cite[Proposition A.6, pp.~41--42]{DE_Hoshi2012}.
We spell out the additional implication needed here,
rather than interpreting that proposition as an
unqualified local section conjecture.

In a free pro-five group the centralizer of a
nonzero power of the primitive generator \(x\)
is \(\langle x\rangle\). One can check this
directly by the Magnus embedding in
\(\mathbb Z_5\langle\!\langle X,Y\rangle\!\rangle\):
for nonzero \(b\in\mathbb Z_5\), the series
\((1+X)^b-1=bX+\cdots\) has a compositional inverse
over \(\mathbb Q_5\). Commuting with it therefore
implies commuting with \(X\). Comparing coefficients
of words containing \(Y\) shows that the centralizer
of \(X\) in the noncommutative power-series algebra
is \(\mathbb Q_5[[X]]\): repeatedly move initial
\(X\)'s to the end of a word containing \(Y\);
once it starts with \(Y\), the commutation equation
forces its coefficient to be zero. Evaluating at \(Y=0\)
then identifies a group element in this subalgebra
with the power of \(x\) determined by its first
abelian coordinate. The same argument applies to
the other two primitive cusp generators, after
a free-basis change. A geometric normalizer of a
nonzero subgroup of \(\langle x\rangle\) is its
centralizer: abelianization forces its exponent
automorphism to be the identity.

Consequently, if a local section has nonzero
specialized inertia in a cusp, its entire specialized
image lies in that cusp's decomposition group.
Indeed the image normalizes the inertia subgroup;
a unit tangential arithmetic lift normalizes the
whole cusp group, and the preceding normalizer
calculation forces the remaining geometric factor
to belong to it. Pullback through (D.67.4) puts the
actual local section in the actual pro-five cusp
packet. This conclusion is compatible with an
ordinary local point close to the cusp having the
same prime-to-\(q\) section image; no point/cusp
separation at the wrong coefficient prime is used.

Now apply this to a section whose exact abelian
cocycle is \((a,0)\). If \(a_q\ne0\), its
inertia image is nonzero and its abelian direction
forces the cusp to be zero: the three cusp
directions are the distinct axes and diagonal.
The parameter in that local packet has Kummer
class \(a|_{G_{\mathbb Q_q}}\), since the primitive
first-coordinate inclusion of its inertia induces
an injection on \(H^1\). Thus this local section
is conjugate to the reference cusp section. The
conjugator can be taken in \(H\). To see this
without ignoring a stabilizer, first combine the
geometric conjugator into the packet with the
cusp-inertia conjugator adjusting its Kummer
cocycle. The abelian image of the resulting
conjugator is fixed by \(G_{\mathbb Q_q}\), since
the two actual abelian cocycles agree. Frobenius
acts there by \(q\), and \(q-1\) is a nonzero
element of \(\mathbb Z_5\), so that fixed image
is zero. Therefore every cocycle (D.67.3) satisfies
\[
 a_q\ne0\quad\Longrightarrow\quad
 c|_{G_{\mathbb Q_q}}\text{ is an }H\text{-coboundary}.
                                                        \tag{D.67.5}
\]
If instead \(a_q=0\), Proposition A.6 forces
the actual section inertia to be zero in
\(\Pi_q^{\rm sp}\). Otherwise its nonzero
cuspidal inertia would have a nonzero abelian
image, contrary to the abelian inertia pair
\((0,0)\). The reference cusp section is also
unramified there, and hence
\[
                   a_q=0\quad\Longrightarrow\quad
                           c(I_q)=1.                   \tag{D.67.6}
\]

\paragraph{Uniform capture in actual finite cocycle sets.}
Use the descending iterated Frattini subgroups
\(P_0=P\), \(P_{r+1}=\Phi(P_r)\). They are
characteristic open subgroups with trivial
intersection. Indeed open subgroups of the
finitely generated free pro-five group are
finitely generated, and their Frattini quotients
are finite; every finite five-group has a
terminating Frattini series. Put
\[
       Q_r=P/P_r,\qquad H_r=\operatorname{im}(H\to Q_r).
                                                        \tag{D.67.7}
\]
The exponent of \(Q_r\) divides \(5^r\), since
\(g^{5^r}\in P_r\) for every \(g\in P\).
These quotients induce the actual topology of
\(H\), and are stable under (D.67.2).

Let \(\rho_{a,r}\) be the action on \(H_r\), and set
\[
 S_r(a)=\{5,\infty\}\cup\{q:a_q\notin5^r\mathbb Z_5\},
 \qquad E_r(a)=Z^1(G_{\mathbb Q,S_r(a)},H_r;\rho_{a,r}).
                                                        \tag{D.67.8}
\]
This is a finite set of actual nonabelian
cocycles. For example each cocycle defines a
homomorphism to the finite group
\(H_r\rtimes\operatorname{im}\rho_{a,r}\), with
fixed second coordinate; finite extensions of
bounded degree unramified outside a fixed
finite set form a finite set, by the same
Hermite--Minkowski argument used in D.47.

Both \(S_r(a)\) and the finite action are
constant on every coset of \(5^r{\cal K}\).
For the action this follows from (D.67.2) and
the exponent bound for \(Q_r\). Moreover every
integral cocycle (D.67.3) reduces to an element
of \(E_r(a)\). At a prime outside \(S_r(a)\)
with nonzero \(a_q\), (D.67.5) expresses its
inertia as a conjugation coboundary for
\(\operatorname{Int}(x^{a(\tau)})\). That
automorphism is the identity on \(Q_r\),
because \(a(\tau)\in5^r\mathbb Z_5\).
Thus the reduced cocycle inertia is trivial.
At a prime with \(a_q=0\), use (D.67.6).
The finite action itself has the same support
bound. This proves exact factorization of the
reduced affine representation through the
group in (D.67.8); it is stronger than an
unramified-cohomology-class assertion.

\paragraph{Dense open conditions eliminating nontrivial kernel restrictions.}
Call \(z\in E_r(a)\) exceptional if its
restriction to \(\ker\rho_{a,r}\) is nontrivial.
Define \({\cal V}_r\subset{\cal K}\) by the
following condition for each exceptional \(z\):
there exists a prime \(q\notin S_r(a)\) with
\[
 \rho_{a,r}(\operatorname{Frob}_q)=1,\qquad
 z(\operatorname{Frob}_q)\ne1,\qquad a_q\ne0.
                                                        \tag{D.67.9}
\]
This is open. On a \(5^r{\cal K}\) coset
all finite data are fixed, there are finitely
many exceptional cocycles, and nonzero
coordinates at a finite choice of witness
primes are open conditions.

It is dense by the actual finite affine-image
argument. Given \(a_0+5^N{\cal K}\), choose
\(M\ge\max\{N,r\}\). For every exceptional
\(z\), an element \(\delta\) in the action
kernel with \(z(\delta)\ne1\) gives an actual
element \((z(\delta),1)\) of
\(H_r\rtimes\operatorname{im}\rho_{a_0,r}\).
Chebotarev supplies infinitely many primes
in this conjugacy class. Both tests in
(D.67.9) are invariant under conjugacy, also
for this nonabelian translation group.
Choose distinct witness primes, outside
\(S_r(a_0)\) and outside the finite set where
\(a_{0,q}\notin5^{M+1}\mathbb Z_5\), and add
\(5^M\kappa(q)\) at each. The finite data
are unchanged, while each new witness
coordinate has valuation exactly \(M\).
Thus all tests hold in the prescribed
neighborhood. This uses precisely
\cite[Chapter V, Theorem 3.23]{PBC_MilneCFT}
on finite Galois quotients that have been
shown to contain the desired element.

Let \({\cal V}\) be the open dense set of
parameters with \(a_q\ne0\) at some
\(q\ne5\). Completeness of the parameter
space, proved in D.66, shows that
\[
                   {\cal G}_0={\cal V}\cap\bigcap_{r\ge1}{\cal V}_r
                                                        \tag{D.67.10}
\]
is a dense \(G_\delta\).

\paragraph{Descent to the actual compact action image.}
Fix \(a\in{\cal G}_0\) and a cocycle
\(c\) as in (D.67.3). Its reduction
\(c_r\in E_r(a)\) cannot be exceptional.
At its witness prime, (D.67.5) would give
an \(H_r\)-coboundary, whose value at a
Frobenius with identity action is one,
contradicting (D.67.9). Thus every \(c_r\)
is trivial on \(\ker\rho_{a,r}\).

It follows that \(c\) is trivial on the
kernel of the actual action on \(H\), and
descends continuously to the compact image
\[
          I_a=\operatorname{im}\left(
             G_{\mathbb Q}\longrightarrow
                         \prod_r\operatorname{Aut}(H_r)\right).
                                                        \tag{D.67.11}
\]
Here the compatible finite actions define
the actual action on \(H\); no enlargement
to an abstract automorphism group is made.
The quotient map from the compact Galois
group proves continuity of the descended
cocycle.

Choose \(q\ne5\) with \(a_q\ne0\).
The image of its inertia in \(I_a\) is
\[
 N=\{\operatorname{Int}(x^{a_q t})|_H:t\in\mathbb Z_5\}.
                                                        \tag{D.67.12}
\]
It is normal in \(I_a\), since every
\(\sigma_a(g)\) sends \(x\) to
\(x^{\chi_5(g)}\). Its fixed subgroup
in \(H\) is trivial:
\[
              H^N=H\cap C_P(x^{a_q})
                    =H\cap\langle x\rangle=1.          \tag{D.67.13}
\]
The first exponent need not be a unit.
By (D.67.5) the descended cocycle is an
\(H\)-coboundary on \(N\). Conjugate it
by that single element of \(H\) so that
it becomes identically one on \(N\).
For \(g\in I_a\), \(n\in N\), compare
the two cocycle formulas at
\(ng=g(g^{-1}ng)\). They imply
\(n(c(g))=c(g)\). Equation (D.67.13)
therefore makes every value one. This
proves that (D.67.3) is a global
\(H\)-coboundary, and hence proves
the singleton assertion in (D.67.1).

Intersect \({\cal G}_0\) with the dense
\(G_\delta\) of infinitely supported
parameters and with the fixed-\(b=0\)
dense \(G_\delta\) from D.66. This is
the asserted \({\cal G}_{\rm full}\).
The permutations of the punctures are
defined over \(\mathbb Q\), act by
homeomorphisms on completed packet
parameters, and transport the actual
section fibres. They give the assertion
for the other two packets.

Finally \(P=\varprojlim_n T_n\): every
finite five-group is nilpotent, so the
lower-central quotients detect every
finite quotient of \(P\). A compatible
sequence of section classes in \(T_n\)
has compatible actual representatives.
Inductively conjugate a representative
at stage \(n+1\) by a lift of the
conjugator identifying its projection
with the chosen representative at stage
\(n\). Their inverse limit is a
continuous section. Conversely two
sections conjugate at every stage are
conjugate in \(P\): their nonempty
closed conjugator sets in the compact
group \(P\) are nested. Thus compatible
nilpotent towers in the indicated fibre
are exactly the singleton full pro-five
fibre. This argument deliberately starts
from compatible classes and says nothing
about unrelated choices of finite-stage
lifts.
\end{proof}
